\subsection{Status Message}

The status message block type reports information about the current state of the telemetry system in integer codes,
such as the current flight state or flight events that are detected. The format of the status message payload
is described in figure \ref{format:status-message}.

\begin{figure}[H]
    \centering
    \begin{bytefield}[bitwidth=0.03\linewidth]{16}
        \bitheader{0-15} \\
        \wordbox{1}{Measurement Time} \\
        \bitbox{8}{Status Code} \\
    \end{bytefield}
    \caption{Status message data block format}
    \label{format:status-message}
\end{figure}

\blocktimestampexp

\paragraph{Status Code}

A status code is an integer representing information about the rocket's state, the state of the telemetry system, or other information.
The measurement time describes a time that the status code was valid, such as when a state change or event occurs, or any time afterwards.
The same code can be transmitted multiple times to ensure that a receiver gets has an up to date picture of the telemetry system's
operation. For example, the state of the telemetry system could be transmitted periodically. The possible status codes are described in
\ref{table:status-message-codes}.

\begin{table}[H]
    \centering
    \begin{tabular}{@{}ll@{}}
        \toprule
        Status message              & Value             \\
        \midrule
        Systems nominal             & 0x00              \\
        Telemetry in idle state     & 0x01              \\
        Telemetry in airborne state & 0x02              \\
        Telemetry in landed state   & 0x03              \\
        Rocket ascent detected      & 0x04              \\
        Rocket apogee detected      & 0x05              \\
        Rocket descent detected     & 0x06              \\
        Reserved for future use     & 0x07 through 0xFF \\
        \bottomrule
    \end{tabular}
    \caption{Status message codes}
    \label{table:status-message-codes}
\end{table}
